\documentclass[11pt]{article}

% Change "review" to "final" to generate the final (sometimes called camera-ready) version.
% Change to "preprint" to generate a non-anonymous version with page numbers.
\usepackage[preprint]{acl}

% Standard package includes
\usepackage{times}
\usepackage{latexsym}
\usepackage{caption}
% For proper rendering and hyphenation of words containing Latin characters (including in bib files)
\usepackage[T1]{fontenc}
% For Vietnamese characters
% \usepackage[T5]{fontenc}
% See https://www.latex-project.org/help/documentation/encguide.pdf for other character sets

% This assumes your files are encoded as UTF8
\usepackage[utf8]{inputenc}

% This is not strictly necessary, and may be commented out,
% but it will improve the layout of the manuscript,
% and will typically save some space.
\usepackage{microtype}
\usepackage{booktabs} 
\usepackage{multirow} 
% This is also not strictly necessary, and may be commented out.
% However, it will improve the aesthetics of text in
% the typewriter font.
\usepackage{inconsolata}

%Including images in your LaTeX document requires adding
%additional package(s)
\usepackage{graphicx}
\usepackage{amsmath} 
\usepackage{amsfonts}

% If the title and author information does not fit in the area allocated, uncomment the following
%
%\setlength\titlebox{<dim>}
%
% and set <dim> to something 5cm or larger.

\title{Discriminative vs. Generative Approaches for Visual Word Sense Disambiguation: An Empirical Study}

% Author information can be set in various styles:
% For several authors from the same institution:
% \author{Rui Zhou \and Jiawei Pei \and Author n \\
%         Address line \\ ... \\ Address line}
% if the names do not fit well on one line use
%         Author 1 \\ {\bf Author 2} \\ ... \\ {\bf Author n} \\
% For authors from different institutions:
% \author{Author 1 \\ Address line \\  ... \\ Address line
%         \And  ... \And
%         Author n \\ Address line \\ ... \\ Address line}
% To start a separate ``row'' of authors use \AND, as in
% \author{Author 1 \\ Address line \\  ... \\ Address line
%         \AND
%         Author 2 \\ Address line \\ ... \\ Address line \And
%         Author 3 \\ Address line \\ ... \\ Address line}

\author{Rui Zhou \\
  University of Zurich\\
  % Affiliation / Address line 2 \\
  % Affiliation / Address line 3 \\
  \texttt{rui.zhou2@uzh.ch} \\\And
  Jiawei Pei \\
  University of Zurich\\
  % Affiliation / Address line 2 \\
  % Affiliation / Address line 3 \\
  \texttt{jiawei.pei@uzh.ch} \\\And
  Nilaksan Selliah \\
  University of Zurich\\
  % Affiliation / Address line 2 \\
  % Affiliation / Address line 3 \\
  \texttt{nilaksan.selliah@uzh.ch} \\} 

\begin{document}
\maketitle
\begin{abstract}
Visual Word Sense Disambiguation (VWSD) requires selecting the correct image for an ambiguous word given minimal context---a surprisingly difficult task where even large vision-language models struggle. This paper presents a systematic comparison of discriminative embedding methods versus generative reasoning approaches on the SemEval-2023 Task 1 benchmark. We evaluate SigLIP2 (a retrieval-optimized dual-encoder) against Qwen3-VL (an 8B-parameter VLM) across five inference strategies: direct matching, Chain-of-Thought reasoning, sense definition prompting, embedding extraction, and caption-based text similarity. Our experiments yield seven key findings: (1) \textbf{Training objectives trump model scale}---SigLIP2 achieves 72.79\% Hit@1 while Qwen3-VL embeddings yield only 9.07\%, a 63-point gap despite Qwen3-VL being 8$\times$ larger; (2) \textbf{Reasoning helps but has limits}---Chain-of-Thought prompting lifts Qwen3-VL to 65.44\%, yet discriminative methods remain 7.35 points ahead; (3) \textbf{Self-generated reasoning beats injected knowledge}---providing sense definitions (63.07\%) underperforms CoT, suggesting VLMs benefit more from articulating their own interpretations; (4) \textbf{Knowledge augmentation can backfire}---adding definitions to SigLIP2 queries \textit{hurts} performance by up to 11.45 points, exposing a mismatch between structured definitions and the natural captions these models were trained on; (5) \textbf{Text-mediated pipelines lose visual grounding}---generating captions then computing text similarity via Sentence-BERT (48.60\%) performs worst, demonstrating that direct visual reasoning outperforms ``describe then match'' approaches; (6) \textbf{Cascade reranking achieves the best zero-shot results}---combining SigLIP2 retrieval with VLM reranking of top-5 candidates yields 77.32\% Hit@1, outperforming either approach alone; (7) \textbf{Minimal fine-tuning outperforms aggressive adaptation}---LoRA-adapted CLIP with only 500 samples and VLM-generated text augmentation (captions, definitions, paraphrases) achieves 75.59\% Hit@1, while larger training regimes cause overfitting. These findings offer actionable guidance: for best accuracy, use cascade pipelines combining fast retrieval with focused generative reasoning; for resource-constrained settings, use conservative LoRA fine-tuning with diverse text augmentation.
\end{abstract}

\section{Introduction}

Consider the phrase ``bank deposit'': should an image retrieval system return a riverbank or a financial institution? Humans resolve such ambiguities effortlessly, but this task---Visual Word Sense Disambiguation (VWSD)---remains challenging for modern vision-language models. While textual Word Sense Disambiguation has been extensively studied \cite{bevilacqua-etal-2021-recent}, its visual counterpart requires reasoning across both modalities with minimal context.

The SemEval-2023 Task 1 \cite{raganato2023semeval} formalized VWSD as a benchmark with three core challenges: (1) \textbf{sparse context}---typically just one or two words accompany the target; (2) \textbf{fine-grained distractors}---candidates include both alternative word senses and visually similar images from the same domain; and (3) \textbf{zero-shot transfer}---systems must generalize without task-specific training. Despite 96 submissions, zero-shot CLIP achieved only 37.20\% accuracy, underscoring the difficulty.

Two paradigms have emerged for multimodal understanding: \textit{discriminative} dual-encoders optimized for retrieval, and \textit{generative} vision-language models (VLMs) capable of reasoning. Which approach better handles VWSD's unique challenges? This question motivates our study.

We systematically compare:
\begin{itemize}
    \item \textbf{SigLIP2} \cite{zhai2023siglip}: A dual-encoder using sigmoid loss instead of contrastive loss, yielding better-calibrated similarity scores for ranking.
    \item \textbf{Qwen3-VL} \cite{yang2025qwen3}: An 8B-parameter autoregressive VLM evaluated across five inference strategies: direct matching, Chain-of-Thought reasoning, sense definition prompting, embedding extraction, and caption-based text similarity.
    \item \textbf{Cascade Reranking}: A two-stage pipeline combining SigLIP2 retrieval with VLM reranking of top-$K$ candidates.
    \item \textbf{LoRA Fine-tuning}: Parameter-efficient adaptation of CLIP with VLM-generated text augmentation (captions, definitions, paraphrases).
\end{itemize}

\noindent\textbf{Key Contributions:}
\begin{enumerate}
    \item We demonstrate that \textbf{pre-training objectives matter more than scale}: SigLIP2 (1B params) outperforms Qwen3-VL embeddings (8B params) by 63 points, establishing that retrieval-specific training is essential for VWSD.
    \item We show that \textbf{explicit reasoning partially compensates} for retrieval-incompatible representations: CoT prompting improves Qwen3-VL by 3.67 points but still trails discriminative methods.
    \item We provide evidence that \textbf{self-generated reasoning outperforms injected knowledge}: sense definitions hurt both paradigms, while text-mediated caption pipelines lose critical visual grounding.
    \item We demonstrate that \textbf{cascade reranking achieves synergistic improvements}: combining fast retrieval with focused VLM reasoning yields 77.32\% Hit@1, outperforming either approach alone.
    \item We offer \textbf{practical design guidance}: use cascade pipelines for best accuracy; for resource-constrained settings, use conservative LoRA fine-tuning with diverse text augmentation.
\end{enumerate}

The remainder of this paper is structured as follows: Section~\ref{sec:related} surveys prior work on visual disambiguation and multimodal retrieval; Section~\ref{sec:method} details our discriminative and generative approaches; Section~\ref{sec:experiments} presents experimental setup, results, and comparative analysis; Section~\ref{sec:limitations} discusses limitations; and Section~\ref{sec:conclusion} concludes with future directions.

\section{Related Work}
\label{sec:related}

\subsection{Visual Word Sense Disambiguation}

Visual Word Sense Disambiguation was formalized as SemEval-2023 Task 1 by \citet{raganato2023semeval}. The task received 96 submissions from 56 teams across three languages, revealing the complexity of multimodal disambiguation. The zero-shot CLIP baseline achieved only 37.20\% Hit@1, substantially lower than traditional textual WSD benchmarks, emphasizing the need for context-aware cross-modal reasoning.

\textbf{Knowledge Retrieval and Ensemble Approaches.} The top-performing systems employed sophisticated retrieval-augmented pipelines. \citet{zhang2023srcb} (SRCB, 1st place in English with 84.02\% Hit@1) leveraged external knowledge bases including WordNet, BabelNet, Wikipedia, and vocabulary.com to construct comprehensive sense inventories. They retrieved definitions and synonyms for the target word, then collected relevant images from the LAION-5B dataset using CLIP-based retrieval. A SimCSE-based biencoder selected the most suitable definition, which was then matched against candidate images using a large CLIP model trained on LAION-2B. Notably, SRCB translated senses from English for Italian and Farsi, achieving strong cross-lingual transfer.

\textbf{Fine-Grained Contrastive Learning.} \citet{yang2023fcll} (FCLL, 1st place overall with 72.56\% average Hit@1) proposed a Fine-grained Contrastive Language-Image Learning framework that learns hierarchical word-sense-image relationships. Their approach introduced a ``sense auto-complementing'' module to enrich minimal context by establishing relationships between concepts and sentences. They constructed a multilingual multimodal knowledge base for ambiguous words, enabling robust cross-lingual performance. This method demonstrated the importance of explicitly modeling sense-level distinctions rather than treating disambiguation as simple image classification.

\textbf{Hybrid Knowledge-Based Systems.} \citet{dadas2023opi} (3rd place overall, 1st in Farsi with 64.0\% Hit@1) combined multimodal embeddings with knowledge-based approaches. Their system integrated CLIP embeddings with additional information retrieved from Wikipedia and lexical databases, using a learning-to-rank (LTR) model to combine outputs from multiple modules. This hybrid approach demonstrates that structured knowledge can complement learned representations, particularly for languages with limited training data.

Other notable approaches include ensemble methods combining multiple CLIP variants with translation models \cite{patil2023rahul}, text-to-image generation for synthetic data augmentation \cite{li2023augmenters}, and prompt augmentation strategies \cite{ghahroodi-etal-2023-sut}. However, nearly all top systems relied on zero-shot or few-shot approaches due to the limited size of training data (12,869 instances), highlighting the importance of strong pre-training for this task.

\subsection{Vision-Language Pre-Training}

\textbf{Contrastive Dual-Encoders.} CLIP \cite{radford2021clip} pioneered large-scale contrastive vision-language pre-training, training separate image and text encoders to align embeddings in a shared space using InfoNCE loss. However, CLIP's batch-based contrastive objective can lead to suboptimal calibration for retrieval tasks. SigLIP \cite{zhai2023siglip} addresses this by replacing InfoNCE with sigmoid loss, which operates on image-text pairs independently rather than requiring large batches. This pairwise classification approach improves retrieval performance and enables more efficient training. Our work leverages SigLIP2, the second generation of this architecture, demonstrating its effectiveness for fine-grained disambiguation.

\textbf{Generative Vision-Language Models.} Recent advances in autoregressive vision-language models have enabled explicit reasoning over visual content. Models like GPT-4V \cite{openai2023gpt4v}, Qwen-VL \cite{bai2023qwen}, and LLaVA \cite{liu2023llava} unify visual perception with language generation, allowing for interpretable decision-making through Chain-of-Thought prompting \cite{wei2022chain}. While these models excel at open-ended visual question answering, their effectiveness for structured ranking tasks like VWSD remains underexplored. Our work investigates whether generative reasoning can match or exceed discriminative retrieval for disambiguation.

\textbf{Multimodal Disambiguation Tasks.} Beyond VWSD, related work includes image-text retrieval \cite{li2020oscar}, visual grounding \cite{kazemzadeh2014referit}, and multimodal lexical semantics \cite{calabrese2020evilbert}. BabelPic \cite{calabrese2020babelpic} provided a multimodal dataset for non-concrete concepts, demonstrating the challenges of representing abstract meanings visually. Our work extends this line of research by focusing on the specific challenge of sense-level disambiguation under minimal context, which requires both fine-grained semantic understanding and robust cross-modal alignment.

\section{Methodology}
\label{sec:method}

We investigate two fundamental paradigms for Visual Word Sense Disambiguation: (1) \textit{discriminative} embedding-based retrieval, and (2) \textit{generative} vision-language reasoning. We also explore knowledge augmentation strategies and hybrid architectures that combine both paradigms. All methods rank 10 candidate images given a target word and minimal context.

\subsection{Task Formalization}

Given an ambiguous target word $w$, a textual context $c$ containing $w$, and a set of $n=10$ candidate images $\mathcal{I} = \{I_1, I_2, \ldots, I_n\}$, the task is to produce a ranking $\pi$ such that the gold image $I^*$ appears as high as possible. Each method computes a relevance score $s_i$ for each image $I_i$, yielding the final ranking $\pi = \text{argsort}(\{s_1, \ldots, s_n\})$.

\subsection{Discriminative Retrieval}

Discriminative methods encode text and images into a shared embedding space, ranking by similarity scores.

\subsubsection{SigLIP2}

SigLIP2 \cite{zhai2023siglip} employs a dual-encoder architecture with separate vision and text transformers. Unlike CLIP's \cite{radford2021clip} InfoNCE loss, SigLIP2 uses sigmoid loss for pairwise classification:
\begin{equation}
\mathcal{L} = -\sum_{i,j} \left[ z_{ij} \log \sigma(s_{ij}) + (1-z_{ij}) \log (1 - \sigma(s_{ij})) \right]
\end{equation}
where $s_{ij} = f_v(I_i)^\top f_t(T_j)$ is the similarity score, $z_{ij} \in \{0, 1\}$ indicates pairing, and $\sigma$ is the sigmoid function. This formulation provides better-calibrated scores for ranking tasks.

For inference, we construct a query $q$ using the format ``this is a photo of a $w$. $c$'' and encode it alongside all candidate images, ranking by $s_i = \sigma(f_t(q)^\top f_v(I_i))$.

\subsubsection{VLM Embedding Extraction}

As an alternative discriminative approach, we extract embeddings from Qwen3-VL \cite{yang2025qwen3} without generation. Text queries are encoded through the language model with mean pooling over hidden states; images are processed similarly. Ranking uses cosine similarity: $s_i = \text{cosine}(f_{\text{text}}(q), f_{\text{image}}(I_i))$. This tests whether VLM representations transfer to retrieval despite being optimized for generation.

\subsection{Generative Reasoning}

Generative methods leverage vision-language models to produce explicit reasoning about image-text matches.

\subsubsection{Direct Matching}

Qwen3-VL \cite{yang2025qwen3} is an autoregressive VLM that integrates a vision encoder with a causal language model. For each candidate image, the model rates relevance on a 0--10 scale given the target word and context. Scores are normalized for ranking.

\subsubsection{Chain-of-Thought (CoT)}

We extend direct matching with step-by-step reasoning \cite{wei2022chain}: (1) interpret the word's meaning in context, (2) describe what the image shows, (3) assess the match. The model outputs brief reasoning followed by a numeric rating.

\subsubsection{Caption-to-Text Pipeline}

Rather than having the VLM directly rate images, we test whether generating captions and computing text-only similarity is effective. Qwen3-VL generates a description for each candidate image, and Sentence-BERT \cite{reimers2019sentence} computes cosine similarity between the query and each caption. This decouples visual understanding from similarity computation.

\subsection{Knowledge Augmentation}

We investigate whether enriching queries with external knowledge improves disambiguation.

\subsubsection{Sense Definition Enrichment}

We employ a two-stage pipeline using language models to generate visual sense definitions---short, concrete descriptions of what would appear in an image for the intended meaning. For example, ``river bank'' yields \textit{``Sloped land beside a river or stream.''} These definitions can be injected into queries for both discriminative and generative methods.

For discriminative retrieval (SigLIP2), we test three query variants: (1) appending definitions to the original query, (2) replacing context with definitions, and (3) using definitions only. For generative methods (Qwen3-VL), definitions are provided as ``visual clues'' to guide image rating.

\subsubsection{Text-to-Image Synthesis}

As an alternative knowledge augmentation strategy, we generate synthetic images representing the intended word sense using text-to-image diffusion models. Given a query, we first expand it into semantically enriched variants using an LLM, then generate a synthetic image $\hat{v}$ via an API-based diffusion model.

Retrieval then operates entirely in the visual embedding space:
\[
    s(q, v_i) = \cos\left(f(\hat{v}), f(v_i)\right)
\]
where $f(\cdot)$ is a vision encoder. This approach reduces cross-modal mismatch by comparing images directly, using the generated image as a visual proxy for the textual query.

\subsection{Cascade Reranking}

Motivated by the complementary strengths of discriminative retrieval (fast, well-calibrated) and generative reasoning (interpretable, context-aware), we propose a two-stage cascade pipeline:

\textbf{Stage 1: Fast Retrieval.} SigLIP2 ranks all 10 candidate images using embedding similarity, efficiently filtering obvious mismatches.

\textbf{Stage 2: Focused Reranking.} Qwen3-VL with Chain-of-Thought reasoning reranks only the top-$K$ candidates from Stage 1, concentrating computational resources on cases where fine-grained disambiguation is needed.

The final ranking concatenates the VLM-reranked top-$K$ with the remaining SigLIP2-ranked candidates: $\pi_{\text{final}} = [\pi_{\text{VLM}}^{1:K}] \oplus [\pi_{\text{SigLIP2}}^{K+1:10}]$.

\subsection{Parameter-Efficient Fine-Tuning}

To adapt discriminative models for VWSD while preserving their general visual-linguistic capabilities, we employ Low-Rank Adaptation (LoRA) \cite{hu2022lora}. LoRA introduces trainable rank-decomposition matrices into transformer layers while keeping pre-trained weights frozen:
\begin{equation}
h = W_0 x + \frac{\alpha}{r} BA x
\end{equation}
where $W_0$ is the frozen pre-trained weight, $B \in \mathbb{R}^{d \times r}$ and $A \in \mathbb{R}^{r \times d}$ are trainable low-rank matrices, $r$ is the rank, and $\alpha$ is a scaling factor.

We apply LoRA to attention projection layers (\texttt{q\_proj}, \texttt{k\_proj}, \texttt{v\_proj}, \texttt{out\_proj}) of CLIP's text encoder while freezing the vision encoder entirely. Training uses SigLIP-style contrastive loss with hard negatives (other candidate images for the same word) but without in-batch negatives to avoid false negative issues when different instances share semantically similar correct images.

\subsubsection{VLM-Based Text Augmentation}

To combat overfitting on the limited training data (12,869 instances), we augment text queries using VLMs. Three augmentation types are generated and cached:

\textbf{Captions:} Qwen3-VL generates natural descriptions of gold images (e.g., ``A sloped riverbank covered with green grass''), providing image-grounded alternatives to the original context.

\textbf{Definitions:} Qwen3 text models produce visual sense definitions (e.g., ``Sloped land beside a river''), offering semantically precise but stylistically different queries.

\textbf{Paraphrases:} The context phrase is rephrased into multiple variants (e.g., ``river bank'' $\rightarrow$ ``stream bank'', ``creek bank edge''), increasing lexical diversity.

During training, original queries are replaced with augmented versions with probability $p_{\text{aug}}$, with uniform random selection among available augmentation types.

\section{Experiments}
\label{sec:experiments}

\subsection{Dataset and Evaluation Metrics}

We use the SemEval-2023 Task 1 English test set, comprising 463 instances with manually annotated gold images. Each instance includes a target word, minimal context (1--2 words), and 10 candidate images. The dataset emphasizes challenging disambiguation scenarios: at least one candidate represents an incorrect sense of the target word, and other distractors are sampled from the same semantic domain (e.g., different animals for ``crane habitat'').

We evaluate using three standard ranking metrics:

\textbf{Hit@1 (Accuracy)} measures the proportion of queries where the gold image is ranked first:
\begin{equation}
\text{Hit@1} = \frac{1}{|\mathcal{Q}|} \sum_{i=1}^{|\mathcal{Q}|} \mathbb{1}[\text{rank}_i = 1]
\end{equation}
where $\text{rank}_i$ is the position of the gold image for query $i$.

\textbf{Mean Reciprocal Rank (MRR)} emphasizes the rank of the first correct result:
\begin{equation}
\text{MRR} = \frac{1}{|\mathcal{Q}|} \sum_{i=1}^{|\mathcal{Q}|} \frac{1}{\text{rank}_i}
\end{equation}

\textbf{Normalized Discounted Cumulative Gain (NDCG@k)} accounts for graded relevance and position:
\begin{equation}
\text{NDCG@k} = \frac{\sum_{i=1}^{k} \frac{2^{\text{rel}_i} - 1}{\log_2(i+1)}}{\sum_{i=1}^{|\text{rel}|} \frac{2^{\text{rel}_i} - 1}{\log_2(i+1)}}
\end{equation}

Hit@1 and MRR are the official metrics for VWSD; we additionally report NDCG@5 and NDCG@10 to assess ranking quality at different depths.

\subsection{Implementation Details}

All experiments use Hugging Face Transformers on NVIDIA GPUs. We describe implementations for each approach category below.

\textbf{Discriminative Models.} For \textbf{CLIP}, we use \texttt{openai/clip-vit-base-patch32} with the prompt template ``This is [mask]. [target\_phrase]''. For \textbf{SigLIP2}, we use the \texttt{so400m-patch14-384} variant (1B parameters, 384px resolution) with Flash Attention 2, encoding one query and 10 images per forward pass (2--4 GB VRAM).

\textbf{Generative Models.} \textbf{Qwen3-VL 8B} uses bfloat16 precision with temperature 0.0 for deterministic generation (16--24 GB VRAM). \textbf{Sense definitions} are generated using Qwen3-8B text models with few-shot prompting and cached to avoid regeneration. \textbf{Text-to-image synthesis} uses Gemini's Imagen API. For cascade reranking, we use $K=5$ candidates.

\textbf{Fine-Tuning.} We fine-tune \texttt{laion/CLIP-ViT-L-14} using Low-Rank Adaptation (LoRA), targeting attention projection layers (\texttt{q\_proj}, \texttt{k\_proj}, \texttt{v\_proj}, \texttt{out\_proj}) while freezing the vision encoder. We use $r=16$, $\alpha=32$, dropout 0.1, and optimize with AdamW (lr$=5\cdot 10^{-5}$) using cross-entropy loss.

\subsection{Main Results}

Table~\ref{tab:results} presents our main results on the English test set (463 instances). We organize methods into four categories and analyze performance across Hit@1, MRR, and NDCG@10.

\begin{table}[t]
\centering
\small
\resizebox{\columnwidth}{!}{%
\begin{tabular}{l|ccc}
\toprule
\textbf{Method} & \textbf{Hit@1} & \textbf{MRR} & \textbf{NDCG} \\
\midrule
\multicolumn{4}{l}{\textit{SemEval-2023 Winners}} \\
\midrule
SRCB \cite{zhang2023srcb} & 84.02 & 89.56 & -- \\
FCLL \cite{yang2023fcll} & 80.13 & 87.42 & -- \\
OPI \cite{dadas2023opi} & 77.97 & 85.88 & -- \\
\midrule
\multicolumn{4}{l}{\textit{Baseline}} \\
\midrule
CLIP (ViT-B/32) & 61.34 & 74.66 & 80.83 \\
\midrule
\multicolumn{4}{l}{\textit{SigLIP2 (Discriminative)}} \\
\midrule
Base (patch16-512) & 68.03 & 79.81 & 84.77 \\
so400m (patch14-384) & 72.79 & 82.76 & 87.00 \\
\quad + def (append) & 72.57 & 82.05 & 86.42 \\
\quad + def (replace context) & 69.11 & 79.58 & 84.53 \\
\quad + def (only) & 61.34 & 74.46 & 80.67 \\
\midrule
\multicolumn{4}{l}{\textit{Qwen3-VL 8B (Generative)}} \\
\midrule
Embedding & 9.07 & 29.31 & 45.51 \\
Direct Matching & 61.77 & 75.93 & 81.83 \\
\quad + CoT & 65.44 & 77.98 & 83.36 \\
\quad + Sense Def. & 63.07 & 74.89 & 80.88 \\
Caption $\rightarrow$ SBERT & 48.60 & 65.10 & 73.52 \\
\midrule
\multicolumn{4}{l}{\textit{Cascade (Ours)}} \\
\midrule
SigLIP2 + VLM Rerank & 77.32 & 85.75 & 89.24 \\
\midrule
\multicolumn{4}{l}{\textit{Fine-Tuned}} \\
\midrule
CLIP (ViT-L/14) + LoRA & 75.59 & 84.82 & -- \\
\bottomrule
\end{tabular}%
}
\caption{Results on English VWSD test set (463 instances). NDCG is computed at $k$=10. SigLIP2 variants differ in model size: Base (patch16-512) vs. so400m (patch14-384, 1B params). The cascade approach (SigLIP2 retrieval + VLM CoT reranking of top-5) achieves the best zero-shot performance.}
\label{tab:results}
\end{table}

\textbf{Result Summary.} Table~\ref{tab:results} reveals a clear hierarchy across method categories. Among \textit{zero-shot discriminative} methods, SigLIP2 so400m achieves the highest Hit@1 (72.79\%), outperforming zero-shot CLIP by 11.45 points and establishing the effectiveness of sigmoid-based contrastive training for retrieval tasks. Our CLIP baseline (61.34\%) substantially exceeds the official SemEval-2023 baseline (37.20\%) due to improved prompt templates. Within \textit{generative} methods, Chain-of-Thought reasoning (65.44\%) leads, followed by sense definition prompting (63.07\%) and direct matching (61.77\%); notably, the embedding extraction method fails catastrophically (9.07\%), confirming that VLM representations are unsuited for direct similarity computation. The \textit{cascade} approach combining SigLIP2 retrieval with VLM reranking achieves our best zero-shot result (77.32\% Hit@1, 85.75\% MRR), demonstrating that discriminative and generative methods are complementary. For \textit{fine-tuned} methods, LoRA-adapted CLIP (75.59\%) outperforms all zero-shot methods except our cascade; we examined three training regimes and found that minimal adaptation (500 samples, 1 epoch) with text augmentation outperformed larger configurations that exhibited overfitting.

\subsection{Analysis and Discussion}

Our experiments reveal several key findings about discriminative vs. generative approaches for VWSD. Rather than merely reporting performance gaps, we analyze the underlying mechanisms that explain these results.

\textbf{Finding 1: The representation-objective alignment hypothesis.} The 63-point gap between Qwen3-VL embeddings (9.07\%) and SigLIP2 (72.79\%) is not simply about ``retrieval-specific training''---it reveals a deeper principle about how optimization objectives shape representation geometry. SigLIP2's sigmoid loss creates embeddings where cosine similarity directly corresponds to match probability, producing a metric space suitable for ranking. Qwen3-VL's next-token prediction objective, by contrast, optimizes for conditional generation: representations encode ``what should come next'' rather than ``how similar are these inputs.'' The catastrophic failure of VLM embeddings for retrieval suggests that generative pre-training creates fundamentally non-metric representation spaces---a finding with implications for any attempt to repurpose generative models for similarity tasks.

\textbf{Finding 2: Generation as a representation bridge.} When Qwen3-VL generates ratings rather than embeddings, performance jumps to 61.77\%---still below SigLIP2 but dramatically improved. This reveals that the generation process itself serves as a ``bridge'' that converts non-metric representations into usable similarity judgments. The model's language generation capability effectively re-routes the comparison through a different computational path that bypasses the limitations of its embedding space. Chain-of-Thought reasoning adds 3.67 points by making this bridging explicit: articulating the intended word sense creates an intermediate representation that the model can more reliably compare against visual content.

\textbf{Finding 3: The externalization paradox.} Sense definitions (63.07\%) underperform CoT (65.44\%) despite providing semantically equivalent disambiguation signals. This reveals what we term the ``externalization paradox'': information generated by the model during reasoning is more useful than equivalent information provided externally. We hypothesize two mechanisms: (1) \textit{representational alignment}---self-generated interpretations use the model's own conceptual vocabulary, while external definitions may emphasize features the model's vision encoder cannot reliably detect; (2) \textit{attention allocation}---during CoT, the model attends to image regions relevant to its own reasoning chain, while injected definitions may direct attention to irrelevant areas.

\textbf{Finding 4: Distribution shift as the cost of enrichment.} Definition enrichment degrades SigLIP2 by up to 11.45 points, revealing a general principle: knowledge augmentation can introduce harmful distribution shift. Dual-encoders learn implicit associations between linguistic patterns and visual features during pre-training on natural captions. Structured definitions break these associations: ``long-necked wading bird with gray plumage'' activates different neural pathways than ``a crane in wetland,'' even when semantically equivalent. Error analysis reveals three failure modes: semantic drift (definitions mentioning related but incorrect concepts), over-specification (emphasizing absent visual details), and context destruction (discarding the dataset's minimal but crucial context). This finding challenges the common assumption that more information is always beneficial.

\textbf{Finding 5: The visual grounding bottleneck.} Caption-based similarity (48.60\%) performs worst, exposing a fundamental limitation of text-mediated visual reasoning. When the VLM generates a caption, it must compress a 384$\times$384 image into 2--3 sentences---an information bottleneck that inevitably discards disambiguation-relevant details. Sentence-BERT then operates in a semantic space misaligned with visual content, compounding the error. This ``describe then match'' pipeline fails because visual disambiguation often depends on subtle features (object orientation, background context, spatial relationships) that resist verbalization. Direct visual grounding---whether through embeddings or explicit ratings---preserves access to these features throughout the comparison process.

\textbf{Finding 6: Complementarity through error orthogonality.} The cascade pipeline achieves 77.32\% Hit@1, outperforming both components alone. This synergy arises from error orthogonality: SigLIP2 and Qwen3-VL fail on different instances. SigLIP2 struggles with cases dominated by pre-training frequency biases (e.g., ``bank'' strongly associated with financial institutions), while Qwen3-VL's explicit reasoning can override such biases. Conversely, Qwen3-VL's uncalibrated scoring fails on cases requiring fine-grained visual discrimination, where SigLIP2's retrieval-optimized embeddings excel. By limiting VLM reasoning to the top-5 candidates, the cascade concentrates expensive computation on genuinely ambiguous cases while preserving the discriminative model's advantage on clear-cut instances.

\textbf{Finding 7: The preservation-adaptation trade-off.} LoRA fine-tuning achieves 75.59\% Hit@1 with only 500 samples and 1 epoch---larger regimes cause performance degradation. This reveals that VWSD success depends heavily on pre-trained representations: the general visual-linguistic associations learned from web-scale data are more valuable than task-specific patterns learned from 12,869 instances. Aggressive adaptation destroys these associations, causing embedding drift. Text augmentation (paraphrasing, synonym substitution) outperforms additional training data because it increases query diversity without overfitting to VWSD-specific patterns. This finding suggests a general principle for low-data multimodal tasks: preserve pre-trained representations and augment inputs rather than adapting weights.

\textbf{Gap to State-of-the-Art.} Our cascade (77.32\%) narrows the gap to SRCB (84.02\%) to 6.70 points without external knowledge bases. SRCB's advantage stems from retrieving additional training signals from LAION-5B and structured knowledge from WordNet/BabelNet. This suggests that further gains require either (1) scaling inference-time retrieval to access relevant visual examples, or (2) improving pre-training to encode richer sense-level distinctions. Pure model architecture improvements appear to have diminishing returns at this performance level.

\section{Limitations}
\label{sec:limitations}

\textbf{Scope of evaluation.} We evaluate on English only (463 test instances), though the SemEval-2023 task includes Italian and Farsi. Our findings about representation geometry and the externalization paradox should generalize across languages, but the specific performance gaps may differ for languages with less pre-training coverage. The small test set also limits statistical power for fine-grained subgroup analysis.

\textbf{Model scale constraints.} We evaluate only the 8B Qwen3-VL variant due to GPU memory limitations. Larger models (72B+) may exhibit stronger reasoning that could narrow the discriminative-generative gap. However, our core finding---that generation acts as a representation bridge---would likely hold regardless of scale, as it reflects architectural rather than capacity limitations.

\textbf{Unexplored design space.} Several promising directions remain unexplored: (1) fine-tuning SigLIP2 rather than CLIP, which may preserve better retrieval properties; (2) prompt ensembles that combine multiple reasoning strategies; (3) adaptive cascade thresholds that vary $K$ based on first-stage confidence. Our LoRA experiments used batch size 1 on CPU, which may have limited training dynamics.

\textbf{Causal attribution.} While we propose mechanistic explanations (representation geometry, externalization paradox, error orthogonality), we do not provide causal verification through controlled interventions. Our hypotheses are consistent with observed results but alternative explanations may exist.

\section{Conclusion}
\label{sec:conclusion}

This study reveals that Visual Word Sense Disambiguation exposes fundamental differences in how discriminative and generative models represent multimodal information---differences obscured by their similar performance on other vision-language tasks.

\textbf{Three principles for multimodal retrieval.} Our findings converge on three general principles: (1) \textit{Representation geometry is task-specific}---embeddings optimized for generation create non-metric spaces unsuitable for similarity-based retrieval, regardless of model scale; (2) \textit{Generation can bridge representation gaps}---the act of producing text transforms internal representations into a format amenable to comparison, explaining why generative scoring outperforms embedding extraction; (3) \textit{Information must match training distribution}---both knowledge augmentation and text-mediated pipelines fail because they introduce signals misaligned with how models learned to process inputs.

\textbf{The cascade as architectural insight.} Our cascade pipeline (77.32\%) achieves its gains not merely through ``combining approaches'' but through error orthogonality: discriminative models fail on frequency-biased cases while generative models fail on fine-grained discrimination. This suggests a design principle for multimodal systems: identify complementary failure modes and route instances accordingly.

\textbf{Implications beyond VWSD.} These findings have broader relevance: (1) \textit{For retrieval system design}---do not assume larger VLMs can replace retrieval-optimized encoders; use each for its strengths; (2) \textit{For knowledge augmentation}---test whether ``enriching'' inputs actually helps or introduces distribution shift; (3) \textit{For fine-tuning}---preserve pre-trained representations through conservative adaptation rather than aggressive task-specific training.

\textbf{Open questions.} Key unresolved issues include: Can pre-training objectives be designed to produce representations suitable for both generation and retrieval? Can cascade routing be learned rather than fixed? What explains the externalization paradox at a mechanistic level? Addressing these questions would advance not only VWSD but multimodal AI systems generally.

\section*{Author Contributions}

\textbf{Rui Zhou} designed and implemented the complete experimental infrastructure, including:
\begin{itemize}
    \item \textit{SigLIP2 pipeline}: Model loader with support for 10+ variants, embedding-based inference engine with three prompt template strategies (default, with-definition, definition-only), and memory-optimized batch processing.
    \item \textit{Qwen3-VL pipeline}: VLM loader with quantization support (bfloat16/AWQ-INT4), and five distinct inference methods---matching, matching with Chain-of-Thought, definition-based description, direct embedding extraction, and caption-to-text similarity via Sentence-BERT.
    \item \textit{Cascade reranking}: Two-stage pipeline combining SigLIP2 retrieval with VLM reranking, including lazy model loading, configurable top-$K$ selection, and compute optimization (70--90\% reduction).
    \item \textit{Definition generation}: Qwen3 text model integration (4B/8B/14B) with persistent JSON caching, few-shot prompting for visual definitions, and batch generation support.
    \item \textit{LoRA fine-tuning}: Parameter-efficient adaptation framework with SigLIP-style contrastive loss, hard negative mining, and VLM-based text augmentation (captions, definitions, paraphrases).
    \item \textit{Evaluation infrastructure}: MRR, Hit@1, and NDCG metric computation with multi-language support and cumulative result tracking.
\end{itemize}
For the paper, Rui wrote the abstract, introduction, related work (Section~\ref{sec:related}), methodology (Section~\ref{sec:method}), analysis and discussion (Section~\ref{sec:experiments}), limitations (Section~\ref{sec:limitations}), conclusion (Section~\ref{sec:conclusion}), and appendix examples.

\textbf{Jiawei Pei} contributed to [TODO: Add Jiawei's contributions here].

\textbf{Nilaksan Selliah} contributed to [TODO: Add Nilaksan's contributions here].

\section*{Acknowledgments}

We thank the SemEval-2023 Task 1 organizers for providing the VWSD benchmark and the open-source community for releasing pre-trained models. This work was conducted as part of coursework at the University of Zurich.

% Bibliography entries for the entire Anthology, followed by custom entries
%\bibliography{anthology,custom}
% Custom bibliography entries only
\bibliography{custom}

\appendix

\section{Implementation Details and Examples}
\label{sec:appendix}

\subsection{Prompt Templates}

\subsubsection{SigLIP2 Query Formats}

For SigLIP2, we tested four query templates. All queries are lowercased to match pre-training. Given target word $w$=``crane'', context $c$=``crane habitat'', and definition $d$=``Long-necked wading bird in wetland'':

\begin{itemize}
    \item \textbf{default} (baseline): \\ \texttt{``this is a photo of a \{w\}. \{c\}''} \\ $\rightarrow$ ``this is a photo of a crane. crane habitat''

    \item \textbf{with\_definition}: \\ \texttt{``this is a photo of a \{w\}. \{c\}. visual clue: \{d\}''} \\ $\rightarrow$ ``this is a photo of a crane. crane habitat. visual clue: long-necked wading bird in wetland''

    \item \textbf{target\_definition}: \\ \texttt{``this is a photo of a \{w\}. \{d\}''} \\ $\rightarrow$ ``this is a photo of a crane. long-necked wading bird in wetland''

    \item \textbf{definition\_only}: \\ \texttt{``this is a photo of \{d\}''} \\ $\rightarrow$ ``this is a photo of long-necked wading bird in wetland''
\end{itemize}

\subsubsection{Qwen3-VL Generative Prompts}

For generative methods, each candidate image is evaluated independently:

\textbf{Direct Matching:}
\begin{quote}
\small
\texttt{Target word: ``\{w\}'' in context ``\{c\}''. Rate how well this image matches the intended meaning (0-10). Output only the number.}
\end{quote}

\textbf{Chain-of-Thought:}
\begin{quote}
\small
\texttt{Target word: ``\{w\}'' in context ``\{c\}''. Step 1: What meaning of ``\{w\}'' is intended? Step 2: What does this image show? Step 3: How well does it match? Provide brief reasoning, then ``Rating: X'' (0-10).}
\end{quote}

\textbf{Sense Definition:}
\begin{quote}
\small
\texttt{Target word: ``\{w\}'' in context ``\{c\}''. Visual clue: \{d\}. Rate how well this image matches (0-10). Output only the number.}
\end{quote}

\subsubsection{Definition Generation Prompt}

Definitions are generated using Qwen3-8B with few-shot prompting:
\begin{quote}
\small
\texttt{You are a visual lexical semantic expert. Provide SHORT, VISUAL, CONCRETE definitions (5-10 words) describing what you would SEE in an image.}

\texttt{Target word: \{w\}}\\
\texttt{Context phrase: \{c\}}\\
\texttt{Visual definition:}
\end{quote}

Examples from few-shot prompt: ``bank'' + ``river bank'' $\rightarrow$ ``Sloped land beside a river or stream''; ``bass'' + ``bass guitar'' $\rightarrow$ ``Large, low-pitched electric guitar with four strings''.

\subsection{Successful Disambiguation Cases}

\textbf{Example 1: ``crane bird''}
\begin{itemize}
    \item \textbf{Gold image:} A bird (crane) in a wetland environment
    \item \textbf{SigLIP2 ranking:} 1st (correct)
    \item \textbf{Qwen3-VL CoT:} ``I observe a long-necked bird with gray plumage standing in shallow water surrounded by reeds. The word 'habitat' strongly suggests a natural living environment rather than a construction site. This matches the biological meaning of crane. Rating: 10/10''
    \item \textbf{Analysis:} The contextual clue ``habitat'' effectively disambiguates toward the biological sense. SigLIP2 correctly prioritizes the bird image, and CoT reasoning explicitly identifies the key semantic relationship.
\end{itemize}

\textbf{Example 2: ``music keyboard''}
\begin{itemize}
    \item \textbf{Gold image:} A piano/synthesizer keyboard
    \item \textbf{SigLIP2 ranking:} 1st (correct)
    \item \textbf{Qwen3-VL CoT:} ``The image shows a musical instrument with black and white keys arranged in the characteristic pattern of a piano. The context 'music' clearly indicates this is a musical keyboard, not a computer keyboard. Rating: 10/10''
    \item \textbf{Analysis:} Both methods successfully leverage the ``music'' context to distinguish between computer and musical keyboards.
\end{itemize}

\subsection{Failure Cases}

\textbf{Example 3: ``bank river''}
\begin{itemize}
    \item \textbf{Gold image:} A riverbank (natural shoreline)
    \item \textbf{SigLIP2 ranking:} 3rd (incorrect; ranked financial institution 1st)
    \item \textbf{Qwen3-VL CoT:} ``I see a building with columns that resembles a financial institution. However, the context 'river' suggests a natural feature. This is ambiguous. Rating: 5/10''
    \item \textbf{Analysis:} SigLIP2 fails due to the dominance of ``bank'' as a financial term in pre-training data. Qwen3-VL recognizes the ambiguity but produces a low confidence score rather than clearly preferring the riverbank image. This case highlights the difficulty of overcoming pre-training biases.
\end{itemize}

\textbf{Example 4: ``pitch baseball''}
\begin{itemize}
    \item \textbf{Gold image:} A baseball pitch (playing field)
    \item \textbf{SigLIP2 ranking:} 2nd (incorrect; ranked pitching action 1st)
    \item \textbf{Qwen3-VL CoT:} ``The image shows a player throwing a ball, which is the act of pitching in baseball. This matches the sports context. Rating: 9/10''
    \item \textbf{Analysis:} Both methods misinterpret ``pitch'' as the action rather than the location. This reflects the higher frequency of ``pitch'' as a verb (throwing) vs. noun (field) in typical training data. Additional context would be needed to resolve this ambiguity.
\end{itemize}

\textbf{Example 5: ``ear tympanum''}
\label{sec:failure_visual_mismatch}

Figure~\ref{fig:failure_case_ear} illustrates a representative failure case caused by \textbf{visual modality mismatch} and \textbf{texture bias}.

\textbf{Instance Details:}
\begin{itemize}
    \item \textbf{Target Word / Context:} ``drum'' (in the context of ear anatomy).
    \item \textbf{Method:} Synthesis-Based Retrieval (Text $\rightarrow$ Generated Image $\rightarrow$ Retrieval).
\end{itemize}

\textbf{Visual Analysis:}
As shown in Figure~\ref{fig:failure_case_ear}(a), the system correctly synthesized a high-quality, realistic image of an eardrum (tympanic membrane) based on the textual query. However, the retrieval step failed to select the correct ground truth.

\begin{itemize}
    \item \textbf{The Error (Fig.~\ref{fig:failure_case_ear}(b)):} The model retrieved an image of an open mouth/throat. Although semantically incorrect, this image shares significant \textbf{low-level visual attributes} with the generated query: pinkish biological tissue, circular framing, strong flash lighting, and depth cues characteristic of endoscopic photography.
    \item \textbf{The Gold Standard (Fig.~\ref{fig:failure_case_ear}(c)):} The correct image is an \textbf{anatomical diagram} (line drawing/illustration). While it perfectly matches the semantic meaning of ``ear anatomy,'' it differs drastically from the query in terms of visual style (textureless, schematic, white background).
\end{itemize}

\textbf{Root Cause:}
The vision encoder (e.g., CLIP/SigLIP) prioritized \textbf{stylistic and textural similarity} (``fleshy photo'' $\approx$ ``fleshy photo'') over \textbf{semantic correspondence} (``ear photo'' $\neq$ ``ear diagram''). This highlights a critical limitation in cross-modal retrieval: models struggle to bridge the \textbf{domain gap} between realistic imagery (photos) and abstract representations (diagrams, sketches) when the distractor images offer a closer textural match.

\begin{figure}[h]
    \centering
    \begin{minipage}{0.32\textwidth}
        \centering
        \includegraphics[width=\textwidth]{proposalTemplate/images/1_query_7.jpg} % 替换你的文件名
        \caption*{(a) Generated Query}
    \end{minipage}
    \hfill
    \begin{minipage}{0.32\textwidth}
        \centering
        \includegraphics[width=\textwidth]{proposalTemplate/images/2_predicted_image.4466.jpg}
        \caption*{(b) Predicted (Incorrect)}
    \end{minipage}
    \hfill
    \begin{minipage}{0.32\textwidth}
        \centering
        \includegraphics[width=\textwidth]{proposalTemplate/images/3_correct_image.4472.jpg}
        \caption*{(c) Gold (Correct)}
    \end{minipage}
    \caption{A failure case demonstrating \textbf{texture bias}. The model matches the generated eardrum photo (a) with a visually similar throat photo (b) instead of the semantically correct but stylistically different anatomical diagram (c).}
    \label{fig:failure_case_ear}
\end{figure}

\end{document}
